\documentclass[french]{book}
\usepackage{teach}
\usepackage{verbatim}
\usepackage{amsmath,amsfonts,amssymb}
\usepackage{pgfplots}
%\usepackage{geometry}
%\geometry{top=1cm, bottom=1cm, left=2cm, right=2cm}
\usepackage[utf8]{inputenc}
\usepackage[french]{babel}
\redefinecolor{thm}{title}{bgcolor}{alizarin}
\redefinecolor{prop}{title}{bgcolor}{meatbrown}
\redefinecolor{defn}{title}{bgcolor}{olivine}
\definecolor{alizarin}{rgb}{0.82, 0.1, 0.26}
\definecolor{meatbrown}{rgb}{0.9, 0.72, 0.23}
\definecolor{olivine}{rgb}{0.6, 0.73, 0.45}
\usepackage{pgf,tikz,tkz-tab}

\usepackage{mathrsfs}
\usetikzlibrary{arrows}

\begin{document}


\begin{center}
\huge{ Devoir surveillé V}
\end{center}

\section*{Calcul de probabilités}
   On tire au hasard une carte dans un jeu de 32 cartes.
   \begin{enumerate}
      \item Quelle est la probabilité de tirer un trèfle ?
      \item Quelle est la probabilité de tirer une carte noire     ?
      \item Quelle est la probabilité de ne pas tirer un carreau ?
      \item Quelle est la probabilité de tirer une figure (roi, dame ou valet) ?
      \item Quelle est la probabilité de tirer un as ?
      \item Quelle est la probabilité de ne pas tirer un valet noir ? \\
   \end{enumerate}

\section*{Loi de probabilité}
On lance deux dés cubiques équilibrés numérotés de 1 à 6.    On note alors le plus grand des deux numéros sortis.
   \begin{enumerate}
      \item Utiliser un tableau à double entrée pour modéliser la situation.
      \item Quel est l'univers $\Omega$ de toutes les issues possibles ?
      \item Établir la loi de probabilité de l'expérience.
   \end{enumerate}

\section*{Arbre pondéré}

Une urne contient 5 boules indiscernables au toucher : deux bleues et trois rouges.
On dispose également de deux sacs contenant des jetons : l'un est bleu et contient un jeton 
bleu et trois jetons rouges, l’autre est rouge et contient deux jetons bleus et deux jetons rouge. \\

On note l'événement:
\begin{itemize}
\item B :"Tirer une boule bleu"
\item R :"Tirer une boule rouge"
\item b :"Tirer un jeton bleu"
\item r :"Tirer un jeton rouge"\\
\end{itemize}

On extrait une boule de l'urne, puis on tire un jeton dans le sac qui est de la même couleur que la boule tirée.

\begin{enumerate}
\item Combien y a-t-il d’issues possibles ?
\item A l'aide d’un arbre pondéré, quel est la probabilité de chacune de ses issues.
\item Quel est la probabilité de l'événement A : « la boule et le jeton extraits sont de la même
couleur »
\end{enumerate}


\newpage


\begin{center}
\huge{ Devoir surveillé VI}
\end{center}
%\title{Devoir surveillé V}
 
\vspace{1cm}
\pagestyle{empty}

Soit $(O,\overrightarrow{i},\overrightarrow{j})$ un repère orthonormé.

\section*{Application du cours}

\begin{enumerate}

\item Dans ce repère, donner les coordonnées du vecteur $\overrightarrow{u}$.

\item Représenter dans le repère les points A(-3 ;2) et B(-4;3).

\item Déterminer les coordonnées  du vecteur $\overrightarrow{AB}$.

\item Déterminer les coordonnées du vecteur $\overrightarrow{u} + \overrightarrow{AB} $

\item Soit $k \in \mathbb{R}$, déterminer, en fonction de $k$, les coordonnées du vecteur $k\overrightarrow{AB}$. 

\item On pose C(0;3) et D(2,0). Déterminer la distance entre C et D.

\item Soit M le milieu du segment [CD], déterminer les coordonnées du point M.   

\item Soit $\overrightarrow{EF} \left( \begin{array}{c}
-4 \\
6 
\end{array} \right)$
 Montrer que les droites $(EF)$ et $(DC)$ sont colinéaires.

\end{enumerate}

\section*{Problème}

Soit L un point du plan. On note $\mathcal{C}$ le cercle de centre $L$ et de rayon $r$.\\

On rappel que $\mathcal{C}$ est l’ensemble des points qui sont à une distance $r$ du point L. \\

On pose L(3;5) et r=3. On souhaite qu’un point M appartiennent au $\mathcal{C}$. 

Recopier et compléter l’assertion suivante :

$$ M \in \mathcal{C} \Longleftrightarrow \vert \vert  \cdots \vert \vert =  .... $$

On pose M$(x;y)$ avec $x,y \in \mathbb{R}$, calculer alors $\overrightarrow{LM}$.

En déduire alors une équation qui traduit le fait que $M \in \mathcal{C}$.



\newpage

\begin{center}
\huge{ Devoir surveillé VII}
\end{center}
%\title{Devoir surveillé V}
 
\vspace{1cm}
\pagestyle{empty}

\section*{Résolution algébrique d’équation $f(x)=k$ et inéquation $f(x)>k$ }

Résoudre algébriquement l’équation $f(x)=3$ avec : \\

\begin{enumerate}
\item $f(x)=x^2 $
\item $f(x)=\dfrac{1}{x}$
\item $f(x)=\sqrt{x}$
\item $f(x)=-x+ \dfrac{2}{3} $
\end{enumerate}

Résoudre algébriquement l’inéquation $f(x)>2$ avec : \\

\begin{enumerate}
\item $f(x)=x^2 $
\item $f(x)=\dfrac{1}{x}$
\item $f(x)=\sqrt{x}$
\item $f(x)=-x+ \dfrac{2}{3} $
\end{enumerate}

\section*{ Résolution graphique d’équation $f(x)=k$ et inéquation $f(x)>k$ }

On représente la courbe de quatre fonctions notées respectivement $C_f,C_g,C_h$ et $C_k$.

\begin{enumerate}
\item Pour chaque courbe représentative, donner le nom de la fonction que la courbe représente.
\item Dresser le tableau de variation de $f(x)$
\item Résoudre graphiquement $f(x) \leq 1,5$
\item Résoudre graphiquement $k(x) = 2$
\item Résoudre graphiquement $h(x)<g(x)$
\item Résoudre graphiquement $f(x)=h(x)$
\end{enumerate}

\newpage 


\begin{center}
\huge{ Devoir surveillé VIII}
\end{center}
%\title{Devoir surveillé V}

\section*{Série statistique}

\begin{enumerate}
\item Donner une liste de 5 valeurs pour laquelle la médiane vaut 12 et la moyenne vaut 10.

\item Donner une liste de 10 valeurs pour laquelle la médiane est identique au premier quartile.
\item Donner deux séries distinctes comportant chacune 6 valeurs, ayant la même valeur médiane et les mêmes quartiles Q1 et Q3 mais dont l’une a 120 pour valeur maximale, alors que l’autre a 12 pour valeur maximale. \\
\end{enumerate}
 
\section*{Étude de salaire}

On considère la série suivante qui correspond à la répartition des salaires mensuels en euros dans une entreprise.

\begin{center}
   \begin{tabular}{| l || c | c | c | c | c |}
     \hline
     Salaires mensuels & 1 000 & 1 200 & 1 300 & 1 500 &  5000 \\ \hline
     Effectifs & 11 & 9 & 14 & 4 & 1 \\ 
     \hline
   \end{tabular}
 \end{center}

\begin{enumerate}
\item Calculer le salaire moyen et le salaire médian au sein de cette entreprise.

\item Déterminer les quartiles Q1 et Q3 en expliquant la démarche.
Interpréter concrètement la valeur trouvée pour Q3.

\item Pour vous recruter, le directeur de l’entreprise vous présente la valeur du salaire moyen. Est-ce un indicateur pertinent ici ? Expliquer.

\end{enumerate}


\section*{Extractions de données et interprétations}
Dans une maternité, on a relevé les masses \textit{(en kg)} à la naissance de 103 bébés nés en avril 2019.


\begin{center}
   \begin{tabular}{| l || c | c | c | c | c |}
     \hline
     Masse & 2,5 & 2,6 & 3 & 3,1 & 5 \\ \hline
     Effectif & 4 & 3 & 9 & 7 & 1 \\ \hline
	  Effectif cumulés & & & & & \\ \hline
	  Fréquence cumulées & & & & & \\
 
     \hline
   \end{tabular}
 \end{center}

\begin{enumerate}
\item Recopier le tableau et complèter les cases vides.
\item Par simple lecture ou calcul mental, combien de bébés pesaient moins de 3kg ? Plus de 3,1kg ? 
\item Quelle est la fréquence des bébés ne dépassant pas 3,1 kg ?
\item Quel est le pourcentage de bébés pesant plus de 2,5 kg et moins de 3,2 kg ? 
\end{enumerate}

\newpage

\begin{center}
\huge{ Devoir surveillé IX}
\end{center}

\section*{Vecteurs et coefficient directeur}
\begin{enumerate}
\item Soit une droite (AB) tel que A(0;2) et B(1;5). Déterminer un vecteur directeur à la droite (AB) puis le coefficient directeur.
\item Soit une droite (OC) tel que son coefficient directeur soit $3,2$, décrire alors l'équation réduite de la droite.
\end{enumerate}

\section*{Appartenance d'un point à une droite}
Les points suivant appartiennent-ils à la droite $\Delta$ d'équation réduite $y=3x-5$? À la droite $\Delta_2$ ayant une équation cartésienne $4y+4x-2=0$? \textit{ (Pour récapituler vos résultats vous pouvez faire un tableau à 4 ligne et 2 colonne) }

\begin{enumerate}
\item A(-2;11)
\item B($\dfrac{3}{2}$;-1)
\item C(1,67;0)
\item D($\dfrac{7}{2}$;2)
\end{enumerate}

Pouvez vous donner le point d'intersection (si il existe) de la droite $\Delta$ et $\Delta_2$?

\section*{Déterminer une équation de droite}
Déterminer les équations de droite de: 
\begin{enumerate}
\item $\Delta_1$ dirigé par  $\overrightarrow{u} \left( \begin{array}{c} 2 \\ -2 \end{array} \right)$  Et passant par H(4;6)
\end{enumerate}


\section*{Associer une droite à l'équation d'une droite}
\begin{enumerate}
\item Donner une équation de droite est associée à $\Delta$
\item Tracer la droite d'équation réduite $\Delta_1: \; y = x-2$
\item Tracer la droite d'équation cartésienne  $\Delta_2: \; 2y+3x-1=0$
\item Tracer la droite d'équation cartésienne $\Delta_3: \; y-3=0$
\item Tracer la droite d'équation cartésienne $\Delta_4: \; x+5=0$
\end{enumerate}

\section*{Intersection de droite}
\begin{enumerate}
\item Soit deux droites $\Delta: \;  y=2x-1$ et $\Delta': \; y=3x+4$
\item Justifier à l'aide du coefficient directeur que les droites $\Delta$ et $\Delta'$ sont sécantes.
\item Déterminer alors le ou les point d'intersection de $\Delta$ et de $\Delta'$
\item (Bonus) Donner deux équations de droites différentes mais qui représente la même droite.
\end{enumerate}
\end{document}